%==========================================================================
\section{Scope-cut plan}
If development falls behind, we cut features in this order while keeping
every runtime check:
\begin{itemize}
  \item[--] \textbf{Cut Stage 1:} Defer dynamic object construction
    (new\_instance), keeping ordinary construction with new. allocate is not
    cut, since the serializer depends on it.
  \item[--] \textbf{Cut Stage 2:} Restrict invoke to zero-argument methods.
  \item[--] \textbf{Cut Stage 3:} Defer dynamic field mutation (setfield),
    keeping read-only field access (getfield); this also gives up the
    serializer's deserialization.
  \item[--] \textbf{Cut Stage 4:} Defer per-parameter method type queries
    (method\_param\_type), keeping basic metadata queries (class\_name,
    field\_count, method\_name, method\_arity).
\end{itemize}
\textbf{Minimum demonstrable version:} reflect with the metadata queries, plus
type-checked getfield and allocate, on the native target, with the specified runtime aborts
and GC safety (Introspection).

